\chapter{The Fundamental Theorem of Calculus}

This blueprint presents a proof of both parts of the Fundamental Theorem of Calculus (FTC).
We work over $\R$ with the Riemann integral on closed intervals $[a,b]$.

\section{Preliminaries}

\begin{definition}[Continuous function on a closed interval]
  \label{def:continuous_on}
  \uses{}
  Let $f : [a,b] \to \R$. We say $f$ is \emph{continuous on $[a,b]$} if $f$ is continuous
  at every point of $[a,b]$ (using the subspace topology).
\end{definition}

\begin{definition}[Interval integral]
  \label{def:interval_integral}
  \uses{}
  For a function $f : [a,b] \to \R$, the \emph{Riemann integral}
  $\int_a^b f(x)\,dx$ is defined as the common limit of upper and lower Riemann sums
  over partitions of $[a,b]$, when this limit exists. We say $f$ is
  \emph{Riemann-integrable} on $[a,b]$ if this integral exists.
\end{definition}

\begin{definition}[Has derivative at]
  \label{def:has_deriv_at}
  \uses{}
  We say $F : \R \to \R$ \emph{has derivative} $f'$ \emph{at} $x$ if
  \[
    \lim_{h \to 0} \frac{F(x+h) - F(x)}{h} = f'.
  \]
\end{definition}

\begin{definition}[Antiderivative]
  \label{def:antiderivative}
  \uses{def:has_deriv_at}
  A function $F : [a,b] \to \R$ is an \emph{antiderivative} of $f : [a,b] \to \R$ if
  $F'(x) = f(x)$ for every $x \in (a,b)$.
\end{definition}

\section{Key lemmas}

\begin{lemma}[Continuous functions are integrable]
  \label{lem:continuous_integrable}
  \uses{def:continuous_on, def:interval_integral}
  If $f : [a,b] \to \R$ is continuous on the compact interval $[a,b]$, then $f$ is
  Riemann-integrable on $[a,b]$.
\end{lemma}
\begin{proof}
  A continuous function on a compact interval is bounded and uniformly continuous,
  from which one can show the upper and lower Riemann sums converge to a common limit.
\end{proof}

\begin{lemma}[Integral bounds]
  \label{lem:integral_bound}
  \uses{def:interval_integral, lem:continuous_integrable}
  If $f$ is integrable on $[a,b]$ and $|f(x)| \le M$ for all $x \in [a,b]$, then
  \[
    \left|\int_a^b f(x)\,dx\right| \le M \cdot (b - a).
  \]
\end{lemma}
\begin{proof}
  This follows from the monotonicity of the integral:
  $-M \le f(x) \le M$ implies $-M(b-a) \le \int_a^b f \le M(b-a)$.
\end{proof}

\begin{lemma}[Mean Value Theorem]
  \label{lem:mean_value_theorem}
  \uses{def:has_deriv_at, def:continuous_on}
  Let $F : [a,b] \to \R$ be continuous on $[a,b]$ and differentiable on $(a,b)$.
  Then there exists $c \in (a,b)$ such that
  \[
    F(b) - F(a) = F'(c)(b - a).
  \]
\end{lemma}
\begin{proof}
  Apply Rolle's theorem to the auxiliary function
  $g(x) = F(x) - F(a) - \frac{F(b)-F(a)}{b-a}(x-a)$.
\end{proof}

\section{First Fundamental Theorem of Calculus}

\begin{theorem}[First Fundamental Theorem of Calculus]
  \label{thm:ftc1}
  \uses{def:continuous_on, def:interval_integral, def:has_deriv_at, lem:continuous_integrable, lem:integral_bound}
  Let $f : [a,b] \to \R$ be continuous on $[a,b]$. Define $F : [a,b] \to \R$ by
  \[
    F(x) = \int_a^x f(t)\,dt.
  \]
  Then $F$ is differentiable on $(a,b)$ and $F'(x) = f(x)$ for every $x \in (a,b)$.
\end{theorem}
\begin{proof}
  Fix $x \in (a,b)$. We wish to show that
  \[
    \lim_{h \to 0} \frac{F(x+h) - F(x)}{h} = f(x).
  \]
  By the definition of $F$,
  \[
    F(x+h) - F(x) = \int_a^{x+h} f(t)\,dt - \int_a^x f(t)\,dt = \int_x^{x+h} f(t)\,dt.
  \]
  Therefore,
  \[
    \frac{F(x+h) - F(x)}{h} = \frac{1}{h}\int_x^{x+h} f(t)\,dt.
  \]

  Since $f$ is continuous at $x$, for every $\varepsilon > 0$ there exists $\delta > 0$ such that
  \[
    |t - x| < \delta \implies |f(t) - f(x)| < \varepsilon.
  \]
  Now suppose $0 < |h| < \delta$. For all $t$ between $x$ and $x+h$, we have
  $|t - x| \le |h| < \delta$, so $|f(t) - f(x)| < \varepsilon$. Hence
  \begin{align*}
    \left|\frac{F(x+h) - F(x)}{h} - f(x)\right|
    &= \left|\frac{1}{h}\int_x^{x+h} f(t)\,dt - f(x)\right| \\[6pt]
    &= \left|\frac{1}{h}\int_x^{x+h} f(t)\,dt - \frac{1}{h}\int_x^{x+h} f(x)\,dt\right| \\[6pt]
    &= \left|\frac{1}{h}\int_x^{x+h} \bigl(f(t) - f(x)\bigr)\,dt\right| \\[6pt]
    &\le \frac{1}{|h|}\left|\int_x^{x+h} |f(t) - f(x)|\,dt\right| \\[6pt]
    &< \frac{1}{|h|} \cdot |h| \cdot \varepsilon = \varepsilon.
  \end{align*}
  Since $\varepsilon > 0$ was arbitrary, we conclude that
  \[
    F'(x) = \lim_{h \to 0}\frac{F(x+h) - F(x)}{h} = f(x). \qedhere
  \]
\end{proof}

\section{Second Fundamental Theorem of Calculus}

\begin{theorem}[Second Fundamental Theorem of Calculus]
  \label{thm:ftc2}
  \uses{def:interval_integral, def:antiderivative, def:continuous_on, lem:continuous_integrable, lem:mean_value_theorem}
  Let $f : [a,b] \to \R$ be Riemann-integrable, and let $F : [a,b] \to \R$ be a continuous
  function such that $F'(x) = f(x)$ for all $x \in (a,b)$. Then
  \[
    \int_a^b f(x)\,dx = F(b) - F(a).
  \]
\end{theorem}
\begin{proof}
  Let $P = \{x_0, x_1, \ldots, x_n\}$ be any partition of $[a,b]$ with
  $a = x_0 < x_1 < \cdots < x_n = b$. By the Mean Value Theorem
  (Lemma~\ref{lem:mean_value_theorem}), for each subinterval $[x_{i-1}, x_i]$ there
  exists a point $c_i \in (x_{i-1}, x_i)$ such that
  \[
    F(x_i) - F(x_{i-1}) = F'(c_i)(x_i - x_{i-1}) = f(c_i)\,\Delta x_i,
  \]
  where $\Delta x_i = x_i - x_{i-1}$.

  Summing over all subintervals, a telescoping sum gives
  \[
    F(b) - F(a) = \sum_{i=1}^n \bigl(F(x_i) - F(x_{i-1})\bigr) = \sum_{i=1}^n f(c_i)\,\Delta x_i.
  \]
  The right-hand side is a Riemann sum for $f$ on $[a,b]$. Since $f$ is Riemann-integrable,
  as the mesh $\|P\| \to 0$ the Riemann sums converge to $\int_a^b f(x)\,dx$. But the
  left-hand side $F(b) - F(a)$ is a constant independent of the partition, so
  \[
    F(b) - F(a) = \int_a^b f(x)\,dx. \qedhere
  \]
\end{proof}

\section{Remark}

Together, the two parts of the Fundamental Theorem establish that differentiation and
integration are inverse operations: the first part shows that integrating and then
differentiating recovers the original function, while the second part shows that definite
integrals can be evaluated by finding antiderivatives. This result is the cornerstone of
calculus and provides the primary practical method for computing integrals.
